
\sgaStableTarget{sga6.e10.expose}\section*{\underline{EXPOSÉ X}\\[3pt] \underline{\underline{FORMALISM OF INTERSECTIONS ON PROPER ALGEBRAIC SCHEMES}}}
\begin{center}
\underline{by O. Jussila}\\
(with an appendix by A. Grothendieck)
\end{center}

\sgaStableTarget{sga6.e10.sec.1}\section*{1. Compatibility of the filtrations with the composition law}
@@M0@@

\sgaStableTarget{sga6.e10.sec.1.1}\subsection*{1.1.}
First consider a noetherian scheme @@M1@@ and the group @@M2@@ of classes of coherent @@M3@@-modules. For every coherent @@M4@@-module @@M5@@, we shall denote its class in @@M7@@ by @@M6@@.

\sgaStableTarget{sga6.e10.def.1-1-1}\textbf{Definition 1.1.1.} \uline{For every integer} @@M8@@, \uline{let} @@M9@@ \uline{denote the subgroup of} @@M10@@ \uline{generated by the classes} @@M11@@ \uline{of coherent} @@M12@@\uline{-modules} @@M13@@ \uline{such that}
@@M14@@
This defines an increasing filtration @@M15@@ of @@M16@@ such that @@M17@@ for @@M18@@.

\sgaStableTarget{sga6.e10.prop.1-1-2}\textbf{Proposition 1.1.2.} \uline{Let} @@M19@@ \uline{be a noetherian scheme,} @@M20@@ \uline{a coherent} @@M21@@\uline{-module such that} @@M22@@\uline{, and} @@M23@@ (@@M24@@\uline{) the subset} @@M25@@ \uline{of} @@M26@@\uline{. For every} @@M27@@, \uline{let} @@M28@@ \uline{denote the integral closed subscheme associated with} @@M29@@\uline{. Then the length of the} @@M30@@\uline{-module} @@M31@@ \uline{is finite for every} @@M32@@ \uline{and nonzero for only finitely many of the} @@M33@@\uline{, and one has}

